\documentclass{beamer}
\usetheme{Madrid}
\usecolortheme{default}

\usepackage{amsmath}
\usepackage{graphicx}
\usepackage{booktabs}
\usepackage{array}

\title[\textit{De Novo} Protein--Ligand Design]{\textit{De Novo} Protein--Ligand Design Including Protein Flexibility and Conformational Adaptation}
\author{Jakob Agamia \and Martin Zacharias}
\institute{Technical University Munich}
\date{\textit{Bioinformatics}, 2026, 42(2), btag027}

\begin{document}

% ------------------------------------------------------------
% TITLE SLIDE
% ------------------------------------------------------------
\begin{frame}
  \titlepage
\end{frame}

% ------------------------------------------------------------
% OUTLINE
% ------------------------------------------------------------
\begin{frame}{Outline}
  \tableofcontents
\end{frame}

% ============================================================
\section{Introduction}
% ============================================================

\begin{frame}{Introduction}
  \begin{block}{The problem}
    \textit{Designing drugs means finding small molecules that bind to a specific protein --- but proteins are not rigid locks.}
    \begin{itemize}
      \item Many drugs are small organic molecules that bind a target protein and alter its function
      \item The search happens in a virtually infinite chemical space
      \item Traditional methods assume a fixed binding pocket, ignoring induced fit --- a critical limitation
    \end{itemize}
  \end{block}

  \vspace{0.5em}

  \begin{block}{The opportunity}
    \textit{AI structure predictors make protein flexibility available at every evaluation step --- this is the paper's core insight.}
    \begin{itemize}
      \item AlphaFold successors can predict a protein--ligand complex from sequence + SMILES in seconds
      \item Explore this vast chemical space with a stochastic search guided by an AI oracle
    \end{itemize}
  \end{block}
\end{frame}

% ============================================================
\section{Materials and Methods}
% ============================================================

% ------ MC simulation approach ------------------------------
\subsection{MC Simulation Approach}

\begin{frame}{MC Simulation Approach}
  \textit{AI-MCLig grows a ligand from benzene by applying random chemical operations, accepting or rejecting each change based on how much it improves the Chai-1 score.}

  \vspace{0.8em}
  \begin{itemize}
    \item Start from benzene (the simplest aromatic scaffold) as the initial SMILES
    \item At each step: apply one chemical operation from a \textbf{moveset}
    \item Rebuild the full protein--ligand complex with Chai-1 from scratch $\rightarrow$ get new pLDDT score
    \item Accept or reject using the \textbf{Metropolis criterion}:
      \[
        P(\Delta s) = \exp(-\beta \cdot \Delta s)
      \]
  \end{itemize}
\end{frame}

% ------ Movesets --------------------------------------------
\subsection{Movesets}

\begin{frame}{Movesets}
  \begin{columns}[T]
    \begin{column}{0.48\textwidth}
      \begin{block}{9 Atomistic moves \quad $\beta = 50$}
        \textit{Less exploratory --- fine-grained edits.}\\[0.4em]
        Nine atom-level chemical operations (add/remove atoms and groups, change atom or bond types, form/break rings). Probabilities empirically tuned on three protein targets.
      \end{block}
    \end{column}
    \begin{column}{0.48\textwidth}
      \begin{block}{BRICS fragment moves \quad $\beta = 5$}
        \textit{More exploratory --- larger jumps.}\\[0.4em]
        Instead of atom-level edits, the fragment-based variant recombines whole molecular building blocks. Larger score changes per step require a lower $\beta$ to maintain a reasonable acceptance rate.
      \end{block}
    \end{column}
  \end{columns}
\end{frame}

% ------ Scoring ---------------------------------------------
\subsection{Scoring of Compounds During MC Simulation}

\begin{frame}{Scoring of Compounds During MC Simulation}
  \textit{The score is not just Chai-1 --- three drug-property biases steer the search toward synthesisable, soluble, drug-like molecules.}

  \[
    \text{score} = 0.8 \cdot \text{Chai-1} - 0.1 \cdot \text{SA} + 0.05 \cdot \text{ESOL} + 0.05 \cdot \text{QED}
  \]

  \begin{itemize}
    \item \textbf{Chai-1} (0.8): pLDDT confidence of the predicted complex --- proxy for binding affinity
    \item \textbf{SA} ($-0.1$): synthetic accessibility score --- penalises hard-to-synthesise molecules; capped at a minimum complexity threshold so the ligand is large enough to fill the pocket
    \item \textbf{ESOL} ($+0.05$): estimated aqueous solubility (Delaney 2004) --- favours water-soluble compounds
    \item \textbf{QED} ($+0.05$): quantitative estimate of drug-likeness (Bickerton et al.\ 2012, without MW factor) --- favours Lipinski-compliant molecules
  \end{itemize}
\end{frame}

% ------ MD / MMGBSA -----------------------------------------
\subsection{Molecular Dynamics Simulations and MMGBSA Calculations}

\begin{frame}{Molecular Dynamics Simulations and MMGBSA Calculations}
  \begin{columns}[T]
    \begin{column}{0.48\textwidth}
      \begin{block}{What is MMGBSA?}
        \begin{itemize}
          \item Estimates the free energy of binding
          \item Based on MD simulation + implicit solvation
          \item More negative $=$ stronger binding
          \item Typical drug range: $-15$ to $-50$ kcal/mol
        \end{itemize}
      \end{block}
    \end{column}
    \begin{column}{0.48\textwidth}
      \begin{block}{Why use it here?}
        \begin{itemize}
          \item Chai-1 gives a structural confidence score; MMGBSA is classical physics (force field + thermodynamics)
          \item Two methods with different assumptions that \textbf{agree} $\Rightarrow$ stronger evidence of correctness
        \end{itemize}
      \end{block}
    \end{column}
  \end{columns}

  \vspace{0.6em}
  \begin{alertblock}{Result}
    Generated compounds show MMGBSA values between $-19$ and $-57$ kcal/mol --- the same range as experimentally known binders (Table 3).
  \end{alertblock}
\end{frame}

% ============================================================
\section{Results and Discussion}
% ============================================================

% ------ Recovery test ---------------------------------------
\subsection{MC-Based Recovery of Target Ligands}

\begin{frame}{MC-Based Recovery of Target Ligands from Simple Starting Compounds}
  \textit{Does the chemical-space engine work? Replacing Chai-1 with dice similarity to a known ligand, the MC search rediscovers the target compound from benzene in most cases.}

  \vspace{0.6em}
  \begin{itemize}
    \item \textbf{Validation strategy:} swap the oracle score for \textbf{Dice similarity} to a known target compound --- if the search works, it should recover the target
    \item \textbf{Result (Table 1):} 8 out of 11 compounds recovered with Dice score = 1.00; the remaining 3 reached Dice $\geq 0.73$
    \item Some complex compounds got stuck in local minima due to the rough chemical-space landscape --- not a failure, the search still finds chemically close solutions
    \item Raising the temperature (lowering $\beta$) allows more exploration but reduces selectivity --- the two-stage protocol (30 parallel runs $\times$ 5000 steps, then 1 run $\times$ 10\,000 steps) is used instead
  \end{itemize}
\end{frame}

% ------ About these validations -----------------------------

\begin{frame}{About These Validations}
  The two validations are complementary, not redundant:

  \vspace{0.8em}
  \begin{center}
    \renewcommand{\arraystretch}{1.5}
    \begin{tabular}{>{\bfseries}p{3.5cm} p{3.5cm} p{3cm}}
      \toprule
      Validation & What it tests & When \\
      \midrule
      Recovery test (Table 1) & Does the \textbf{search engine} work? & Before real design \\
      \addlinespace
      MMGBSA + Boltz-2 (Table 2) & Are the generated compounds \textbf{plausible binders}? & After real design \\
      \bottomrule
    \end{tabular}
  \end{center}
\end{frame}

% ------ Chai-1 correlation ----------------------------------
\subsection{Correlation of Chai-1 Confidence Score and Ligand Binding Affinity}

\begin{frame}{Chai-1: The Imperfect Oracle}
  \textit{Chai-1 pLDDT is not a direct affinity measurement --- but its average trend is positive, and that is enough to guide a Monte Carlo search uphill.}

  \vspace{0.6em}
  \begin{itemize}
    \item Three protein targets evaluated: bromodomain, Pim-1 kinase, p38 kinase
    \item Experimental binding affinities (kcal/mol) from PDBbind (Liu et al.\ 2017) compared to Chai-1 pLDDT scores
    \item \textbf{Key result:} no direct point-to-point correlation; but the \textbf{average trend} across binned Chai-1 scores is consistently positive for all three targets (Figure 1)
    \item Autodock Vina tested as alternative oracle --- no significant improvement in correlation
    \item \textbf{Conclusion:} the oracle is imperfect but good enough to guide the search
  \end{itemize}
\end{frame}

% ------ De novo generation - results ------------------------
\subsection{\textit{De Novo} Generation Using Atomistic-Step MC Simulation}

\begin{frame}{\textit{De Novo} Generation: Targets and Scores}
  \textit{From benzene to a drug-like compound in 2\,000 steps: AI-MCLig generates novel ligands whose MMGBSA and Boltz-2 scores match experimentally known binders.}

  \vspace{0.6em}
  \begin{itemize}
    \item Applied to four targets: bromodomain, p38 kinase, Pim-1 kinase, $\beta$-1 adrenergic receptor (GPCR)
    \item 2\,000 MC steps, $\beta = 50$, starting from benzene; \textbf{3 independent runs per target} (${\sim}20$ h each on NVIDIA RTX4090)
    \item Final Chai-1 scores: \textbf{0.85--0.92} across all 12 compounds (Table 2)
    \item MMGBSA scores: \textbf{$-19$ to $-57$ kcal/mol} --- in the same range as experimentally validated high-affinity binders (Table 3)
  \end{itemize}
\end{frame}

% ------ De novo generation - convergence --------------------

\begin{frame}{\textit{De Novo} Generation: Convergence and Protein Adaptation}
  \begin{itemize}
    \item Score improves rapidly in the \textbf{first 250 steps}, then more gradually (Figure 3A)
    \item Pocket RMSD (Figure 3C):
      \begin{itemize}
        \item \textbf{Bromodomain} stays low (${\sim}0.5$ \AA) --- rigid pocket
        \item \textbf{p38 kinase} and \textbf{Pim-1 kinase} show larger fluctuations (0.5--2.0 \AA) --- flexible pockets adapting to each candidate ligand
      \end{itemize}
    \item \textbf{Cleanup phase:} 100 additional steps optimising SA, ESOL, QED with a Chai-1 threshold constraint --- improves drug-likeness without losing the binding score
    \item UMAP analysis: structurally diverse outputs per target all land near known-binder clusters in chemical space
  \end{itemize}
\end{frame}

% ============================================================
\section{Conclusion}
% ============================================================

\begin{frame}{Conclusion}
  \textit{AI-MCLig demonstrates that a simple Monte Carlo search scored by an AI structure predictor is sufficient to generate novel drug-like ligands with predicted binding energies comparable to known experimental binders.}

  \vspace{0.6em}
  \begin{itemize}
    \item \textbf{Central claim:} Chai-1 as a scoring oracle inside an MC loop enables \textit{de novo} ligand generation with \textbf{full protein flexibility} --- no rigid pocket assumption
    \item \textbf{Validated on 4 targets:} bromodomain, p38 kinase, Pim-1 kinase, $\beta$-1 adrenergic receptor --- diverse target classes
    \item \textbf{MMGBSA and Boltz-2 scores} of generated compounds overlap with experimentally known high-affinity binders (Tables 2 \& 3)
    \item \textbf{Complementary, not competing:} supplements traditional docking and deep generative models; adaptable to scaffold optimisation, side-chain refinement, or fragment recombination
  \end{itemize}
\end{frame}

% ============================================================
\section*{Q\&A}
% ============================================================

\begin{frame}{Ask Us Something}
  \begin{itemize}
    \item Why Chai-1 as oracle rather than a formula-based physical energy evaluation?
    \item Why those specific probabilities for each move?
    \item What is $\beta$? (Metropolis acceptance criterion)
    \item Difference between \textit{ligand} and \textit{binder}
    \item Why start with benzene?
    \item Is high affinity always the functional requirement?
    \item Once the score plateaus after ${\sim}250$ moves, the MC engine seems inefficient --- how would you fix it?
  \end{itemize}
\end{frame}

\end{document}
